\section*{Práctica de laboratorio: Cinemática galileana}

\section*{Objetivos}

\begin{itemize}
\item Analizar experimentalmente el movimiento de proyectiles en dos dimensiones.
\item Caracterizar experimentalmente un sistema de lanzamiento de proyectiles (mini lanzador).
\item Determinar experimentalmente las variables cinemáticas asociadas al disparo de un proyectil: velocidad inicial, ángulo de lanzamiento y alcance.
\item Estudiar el movimiento acelerado de un objeto que desciende por un plano inclinado.
\item Determinar la aceleración del carro en función del ángulo de inclinación del plano.
\item Integrar dos sistemas físicos independientes (lanzador y carro) para resolver un problema de interceptación entre proyectil y objeto en movimiento.
\item Aplicar herramientas de análisis de video y sensores electrónicos para la obtención de datos experimentales.
\item Diseñar estrategias experimentales de calibración para lograr que un proyectil impacte un objeto móvil.
\end{itemize}

\section*{Conceptos a estudiar}

\begin{itemize}
\item Movimiento parabólico de proyectiles.
\item Descomposición del movimiento en componentes horizontal y vertical.
\item Movimiento uniformemente acelerado.
\item Movimiento en plano inclinado.
\item Velocidad inicial y ángulo de lanzamiento.
\item Alcance de un proyectil.
\item Condiciones de intersección entre dos trayectorias.
\item Calibración experimental de sistemas físicos.
\item Análisis de movimiento mediante video (software Tracker).
\item Medición mediante sensores electrónicos (Arduino).
\end{itemize}

\rule{\textwidth}{1pt}

\section*{Fundamento teórico}

El movimiento de un proyectil lanzado con velocidad inicial $v_0$ y ángulo $\theta$ respecto a la horizontal puede describirse separando el movimiento en dos componentes independientes.

\subsection*{Movimiento horizontal}

\begin{equation}
x(t) = v_0 \cos\theta \; t
\end{equation}

En la dirección horizontal no existe aceleración (si se desprecia la resistencia del aire), por lo que la velocidad es constante.

\subsection*{Movimiento vertical}

\begin{equation}
y(t) = v_0 \sin\theta \; t - \frac{1}{2} g t^2
\end{equation}

donde $g$ es la aceleración gravitacional.

La trayectoria resultante es parabólica.

\subsection*{Movimiento en plano inclinado}

Cuando un objeto desciende por un plano inclinado de ángulo $\alpha$, la componente de la gravedad paralela al plano genera una aceleración:

\begin{equation}
a = g \sin\alpha
\end{equation}

Si el carro parte del reposo, su posición sobre la rampa se describe mediante:

\begin{equation}
x(t) = \frac{1}{2} a t^2
\end{equation}

\textbf{Importante:} Estas ecuaciones describen el sistema cuando el marco de referencia tiene un grado de libertad paralelo a la dirección de movimiento. \textbf{¡¡En esta práctica, NO se tomará ese marco de referencia!!} \textit{Las ecuaciones deben ser modificadas}.

\subsection*{Problema de intercepción}

Cuando un proyectil es disparado hacia un objeto en movimiento, el impacto ocurre si ambos ocupan la misma posición espacial en el mismo instante de tiempo.

Esto implica que:

\begin{equation}
x_{\text{proyectil}}(t) = x_{\text{carro}}(t)
\end{equation}

\begin{equation}
y_{\text{proyectil}}(t) = y_{\text{carro}}(t)
\end{equation}

En esta práctica se caracterizarán ambos sistemas de manera independiente para posteriormente ajustar los parámetros experimentales y lograr que el proyectil impacte el carro mientras este desciende por la rampa.

\rule{\textwidth}{1pt}

\section*{Materiales}

\begin{itemize}
\item Patín de dinámica (carro)
\item Plano inclinado
\item Mini lanzador de proyectiles PASCO
\item Balines metálicos
\item Láminas de fórmica
\item Cinta de enmascarar
\item Metro o cinta métrica
\item Cronómetro
\item Cámara para análisis de video
\item Computador con software \textbf{Tracker}
\item Arduino con sensores de detección
\end{itemize}

Observación importante:

El mini lanzador posee \textbf{tres niveles de potencia}. Debido a posibles variaciones mecánicas, algunos lanzadores pueden dispararse espontáneamente en ciertos niveles. Por esta razón se recomienda trabajar únicamente con \textbf{dos niveles de lanzamiento} como mínimo.

\rule{\textwidth}{1pt}

\section*{Actividad I: Caracterización del movimiento en un plano inclinado mediante videoanálisis}

\textbf{Objetivo:} Determinar experimentalmente la aceleración del carro mediante el ajuste cuadrático de sus ecuaciones de movimiento $x(t)$ y $y(t)$, y analizar la consistencia de los resultados para distintas condiciones iniciales.

\subsection*{Sistema experimental}

Un carro se deja deslizar desde el reposo sobre un plano inclinado. El movimiento es registrado en video para su posterior análisis utilizando software de videoanálisis.

\subsection*{Procedimiento}

\begin{itemize}
\item Ajustar el plano inclinado a un ángulo fijo $\alpha$.
\item Seleccionar tres alturas iniciales distintas: $h_1$, $h_2$ y $h_3$.
\item Para cada altura, grabar al menos \textbf{10 videos} del movimiento del carro.
\item Asegurarse de que:
  \begin{itemize}
  \item La cámara esté fija durante la grabación.
  \item El plano inclinado sea completamente visible.
  \item Exista una referencia de escala (regla o longitud conocida).
  \end{itemize}
\item Importar cada video en el software Tracker.
\item Calibrar la escala y definir un sistema de referencia adecuado (eje $x$ paralelo al plano inclinado).
\item Marcar la posición del carro cuadro a cuadro para obtener los datos de posición en función del tiempo.
\item Exportar los datos $x(t)$ y $y(t)$.
\end{itemize}

\subsection*{Análisis de datos}

Para cada video:

\begin{itemize}
\item Ajustar los datos experimentales a funciones cuadráticas de la forma:

\begin{equation}
x(t) = A_x t^2 + B_x t + C_x
\end{equation}

\begin{equation}
y(t) = A_y t^2 + B_y t + C_y
\end{equation}

\item Registrar los parámetros $A$, $B$ y $C$ obtenidos en cada ajuste.
\item Interpretar físicamente los parámetros:
  \begin{itemize}
  \item $A \propto a$ (relacionado con la aceleración).
  \item $B \propto v_0$ (relacionado con la velocidad inicial).
  \item $C \propto r_0$ (relacionado con la posición inicial).
  \end{itemize}
\item Determinar la aceleración experimental a partir de las ecuaciones de movimiento del sistema.

\item Calcular el promedio y la desviación estándar de la aceleración para cada altura ($h_1$, $h_2$, $h_3$).
\item Comparar los valores obtenidos entre sí.
\item Comparar con el valor teórico.
\end{itemize}

\subsection*{Análisis adicional}

\begin{itemize}
\item Analizar si la aceleración depende de la altura inicial.
\item Discutir posibles fuentes de error:
  \begin{itemize}
  \item Fricción
  \item Precisión en el marcado de puntos
  \item Resolución temporal del video
  \end{itemize}
\item Comparar los valores de $A_y$ con los de $A_x$ y discutir su significado físico.
\end{itemize}

\subsection*{Resultados esperados}

\begin{itemize}
\item Tablas, histogramas o cualquier método de presentación de datos, con los parámetros $A$, $B$, $C$ para cada video.
\item Promedios y desviaciones estándar de la aceleración.
\item Gráficas de $x(t)$ y $y(t)$ y sus ajustes cuadráticos.
\item Comparación entre resultados experimentales y teóricos.
\item Discusión sobre la validez del modelo de movimiento uniformemente acelerado.
\end{itemize}

\rule{\textwidth}{1pt}

\section*{Actividad II: Caracterización del lanzador de proyectiles y del resorte}

\textbf{Objetivo:} Determinar experimentalmente la velocidad de salida de una esfera metálica para distintos niveles de disparo y, a partir de ello, estimar la constante elástica del resorte del lanzador.

\subsection*{Sistema experimental}

El mini lanzador impulsa una esfera metálica mediante un resorte comprimido. La velocidad de salida depende del nivel de disparo seleccionado, el cual determina la compresión inicial del resorte.

\subsection*{Procedimiento}

Se debe escoger una de las siguientes opciones experimentales:

\begin{itemize}
\item \textbf{Opción A:}
  \begin{itemize}
  \item Fijar el lanzador en el nivel de disparo 1 y luego en el nivel 2.
  \item Para cada nivel, realizar \textbf{15 disparos}.
  \end{itemize}

\item \textbf{Opción B:}
  \begin{itemize}
  \item Fijar el lanzador en los niveles de disparo 1, 2 y 3.
  \item Para cada nivel, realizar \textbf{10 disparos}.
  \end{itemize}
\end{itemize}

Para cada disparo:

\begin{itemize}
\item Registrar el tiempo de interrupción de la fotocompuerta utilizando Arduino.
\item Medir la masa de la esfera con una gramera.
\item Medir el diámetro de la esfera utilizando un pie de rey.
\item Medir la compresión del resorte para cada nivel de disparo.
\end{itemize}

\textbf{Opcional (muy opcional):}

\begin{itemize}
\item Grabar el movimiento del proyectil.
\item Analizar el video con el software Tracker para obtener la trayectoria y comparar con los resultados obtenidos por la fotocompuerta.
\end{itemize}

\subsection*{Análisis de datos}

\begin{itemize}
\item Calcular la velocidad de salida de la esfera para cada disparo mediante:
\begin{equation}
v = \frac{d}{\Delta t}
\end{equation}
donde $d$ es el diámetro de la esfera y $\Delta t$ el tiempo de interrupción de la fotocompuerta.

\item Para cada nivel de disparo:
  \begin{itemize}
  \item Calcular la velocidad promedio.
  \item Determinar la desviación estándar.
  \end{itemize}

\item Considerando conservación de energía:
\begin{equation}
\frac{1}{2} k x^2 = \frac{1}{2} m v^2
\end{equation}

\item Despejar la constante elástica del resorte:
\begin{equation}
k = \frac{m v^2}{x^2}
\end{equation}

\item Calcular el valor de $k$ para cada nivel de disparo.

\item Analizar si el valor de $k$ es consistente entre los distintos niveles.
\end{itemize}

\subsection*{Análisis adicional}

\begin{itemize}
\item Discutir posibles fuentes de error:
  \begin{itemize}
  \item Pérdidas por fricción
  \item Energía no transferida completamente a la esfera
  \item Precisión en la medición de $x$, $d$ y $\Delta t$
  \end{itemize}

\item Comparar, si se realizó, la velocidad obtenida mediante Tracker con la calculada por la fotocompuerta.

\item Evaluar la validez del modelo de conservación de energía en el sistema real.
\end{itemize}

\subsection*{Resultados esperados}

\begin{itemize}
\item Tabla de tiempos de interrupción y velocidades calculadas.
\item Velocidades promedio por nivel de disparo.
\item Valores estimados de la constante elástica $k$.
\item Comparación entre distintos niveles de disparo.
\item Discusión sobre la consistencia del modelo físico.
\end{itemize}

\rule{\textwidth}{1pt}

\section*{Actividad III: Impacto sobre objeto en movimiento}

\textbf{Objetivo:} Ajustar experimentalmente los parámetros del sistema para lograr que el proyectil impacte el carro mientras este desciende por la rampa.

\subsection*{Sistema experimental}

Se combinan los dos sistemas estudiados previamente:

\begin{itemize}
\item El lanzador de proyectiles.
\item El carro descendiendo por el plano inclinado.
\end{itemize}

\subsection*{Diseño experimental}

Fijar los siguientes parámetros:

\begin{itemize}
\item El ángulo del lanzador.
\item La posición inicial del carro sobre la rampa.
\end{itemize}

Determinar experimentalmente los valores adecuados de los demás parámetros para lograr la intercepción.

\subsection*{Procedimiento}

\begin{itemize}
\item Colocar el lanzador en posición fija.
\item Seleccionar un nivel de potencia previamente caracterizado.
\item Colocar el carro en una posición inicial sobre la rampa.
\item Liberar el carro y disparar el proyectil de manera sincronizada.
\item Observar si ocurre impacto entre el proyectil y el carro.
\item Ajustar la posición inicial del carro o el momento del disparo.
\item Repetir el procedimiento hasta lograr que el balín impacte el carro.
\end{itemize}

\subsection*{Análisis físico}

Para analizar el impacto se deben considerar:

\begin{itemize}
\item La velocidad inicial del proyectil.
\item El tiempo de vuelo del proyectil.
\item La aceleración del carro.
\item La posición inicial del carro sobre la rampa.
\end{itemize}

\textbf{Estimar:}

\begin{itemize}
\item El tiempo en que el proyectil alcanza la altura del carro.
\item La posición del carro en ese mismo instante.
\end{itemize}

\subsection*{Resultados esperados}

\begin{itemize}
\item Registro de los intentos realizados.
\item Configuración experimental final que produce el impacto.
\item Análisis teórico comparado con la observación experimental.
\end{itemize}

\rule{\textwidth}{1pt}

\section*{Preguntas orientadoras}

\begin{itemize}
\item ¿Qué variables determinan el alcance de un proyectil?
\item ¿Cómo influye el ángulo de lanzamiento en la trayectoria del balín?
\item ¿Por qué el movimiento horizontal y vertical del proyectil pueden analizarse de manera independiente?
\item ¿Cómo se determina experimentalmente la aceleración de un objeto en un plano inclinado?
\item ¿Qué condiciones deben cumplirse para que el proyectil impacte el carro?
\item ¿Qué fuentes de error pueden afectar la precisión del experimento?
\item ¿Cómo podrían mejorarse las mediciones utilizando sensores electrónicos?
\end{itemize}

\rule{\textwidth}{1pt}

\section*{Conclusión esperada}

Al finalizar la práctica, el estudiante deberá comprender cómo caracterizar experimentalmente sistemas físicos independientes y cómo integrar dichos sistemas para resolver un problema dinámico de interceptación entre dos objetos en movimiento. Además, deberá desarrollar habilidades experimentales en medición, análisis de datos y diseño de estrategias de calibración para lograr un objetivo físico específico dentro de un sistema dinámico.